
(i). D'autre part, (iii) implique (ii) gr\^ace \`a 2.3 et 2.8.

\subsection*{4. Th\'eor\`emes de finitude pour $\underline{\Pic}^\tau_{X/S}$}

\textbf{4.1.} Soient \(S\) un sch\'ema et \(G\) un foncteur en groupes commutatifs sur \(S\). On d\'esigne par \(G^0\) (resp. \(G^\tau\)) le sous-foncteur en groupes dont les \(T\)-points \(f:T\to G\) sont ceux de \(G\) tels que pour tout corps alg\'ebriquement clos \(K\) et tout \(K\)-point \(t\) de \(T\), il existe un \(k\)-sch\'ema connexe \(Z\), deux \(K\)-points \(z,z_0\) de \(Z\) et un morphisme \(g:Z\to G\) tels que \(g(z)=f(t)\) et \(g(z_0)=0\) dans \(G(K)\) (resp. il existe un entier non-nul \(n\) tel que \(n.f(t)\in G^0(K)\)).

Il est \'evident que la formation de \(G^0\) (resp. \(G^\tau\)) commute \`a tout changement de base, qu'un morphisme de foncteurs en groupes \(G\to G'\) induit un morphisme \(G^0\to G'^0\) (resp. \(G^\tau\to G'^\tau\)), et que si \(S\) est un corps et \(G\) est repr\'esent\'e par un sch\'ema en groupes localement de type fini, alors \(G^0\) est repr\'esent\'e par la composante connexe de l'\'el\'ement neutre, qui est un sous-sch\'ema ouvert (et ferm\'e), de type fini, et g\'eom\'etriquement irr\'eductible (SGA 3 VI\(_A\) 2.4).

\textbf{Lemme 4.2.} \emph{Soit \(\Phi:G\to G'\) un morphisme entre deux foncteurs en groupes commutatifs sur un sch\'ema \(S\). On suppose que pour tout corps alg\'ebriquement clos \(K\) et tout morphisme \(\Spec(K)\to S\), \(G_K\) (resp. \(G'_K\)) soit repr\'esentable par un sch\'ema en groupes localement de type fini, et \(\Phi_K\) soit repr\'esentable par un morphisme de type fini, (c-\`a-d, quasi-compact). Alors}
\[
  \Phi^{-1}(G'^\tau)=G^\tau.
\]

